% ADD THIS HEADER TO ALL NEW CHAPTER FILES FOR SUBFILES SUPPORT
\documentclass[../CLthesis.tex]{subfiles}
\graphicspath{{\subfix{../images/}}}
% END OF SUBFILES HEADER

\begin{document}
\chapter{Discussion}%
\label{chap:discussion}

The preceding three empirical chapters each answered one facet of the central research question in isolation: who is targeted (Chapter~\ref{chap:who}), how hatred is rhetorically constructed (Chapter~\ref{chap:how}), and what moral motivations drive it (Chapter~\ref{chap:why}). This chapter does the work that results sections cannot do alone --- it explains \emph{why} these patterns emerge, evaluates the computational framework that produced them, and reflects honestly on what the pipeline can and cannot claim. The chapter is organised around three tasks: synthesising the cross-dimensional findings into a coherent account of cross-lingual hate speech structure (Section~\ref{sec:discussion-synthesis}); assessing the scientific rationale of the who--how--why pipeline as a methodological contribution (Section~\ref{sec:discussion-pipeline}); and mapping the sources of error and technical limitation that bound the conclusions that can be drawn (Section~\ref{sec:discussion-errors}).

% ------------------------------------------------------------------
\section{Synthesis of Cross-Dimensional Findings: The Who--How--Why Triad}
\label{sec:discussion-synthesis}
% ------------------------------------------------------------------

The most important single finding of this thesis is not contained in any individual results chapter, but in the relationship between them: \emph{attack frame selection in online hate speech is not arbitrary but is systematically determined by the social category of the target}. This causal link between who is attacked and how they are attacked is the structural regularity that connects RQ1 to RQ2, and it persists across all three language communities despite their divergent discourse styles.

\subsection{How Target Type Determines Rhetorical Frame}

The pattern is clearest for clerical targets. Across English, Chinese, and Japanese corpora alike, role-based targets at the level of the individual clergy member --- English \textit{priests}, Chinese 神父, and by extension 主教~(\textit{bishops}) and 教皇~(\textit{the Pope}) --- attract the Sexual Abuse Frame as their primary attack vector. In Chinese, the frame's dominance is striking: 55\% of all predicates directed at 神父 are abuse-related, and the same frame leads for 主教 (23 predicates) and 教皇 (18). In English, \textit{priests} receive the Sexual Abuse Frame in 50\% of their predicate assignments. This cross-lingual convergence suggests that clerical role-categories activate a culturally shared discourse \emph{genre}: the priest is a specific kind of moral authority figure, and moral authority's betrayal through sexual misconduct constitutes a predictable, genre-appropriate attack vector regardless of the speaker's native language.

The Japanese corpus provides the clearest counterexample and therefore the strongest confirmation of the same principle. The Unification Church (統一教会) is not a clerical role-category but a political-institutional target, and it attracts not the Sexual Abuse Frame but \textit{Political Interference} (31\% of predicates, 80 instances) and \textit{Economic Exploitation} (20\%, 51 instances). The Liberal Democratic Party (自民党), as the political patron enabling the Church, is attacked almost exclusively through the Political Interference frame (67\%). Shinzo Abe (安倍晋三), as a named political individual, receives Moral Accusation alongside Political Interference --- encoding him simultaneously as a morally culpable person and a political actor. The target's perceived social role --- religious-political institution, enabler party, and named individual --- constrains the available rhetorical repertoire in each case. Importantly, \textit{霊感商法}~(spirit-based fraud) vocabulary such as 献金, 詐欺, and 搾取 appears nowhere in the English or Chinese targets, precisely because the financial predation narrative is specific to the social context of the Unification Church scandal rather than to Christianity as a religious category.

For collective identity labels --- English \textit{Christians}/\textit{Catholics}, Chinese 基督徒, Japanese キリスト教徒 --- the dominant frames shift again: Rhetorical Attack and Moral Accusation together account for the majority of predicate assignments across all three languages. Generic membership categories invite generic condemnation; they do not anchor the more specific institutional and misconduct frames that require a specific actor or organisation to accuse. This tripartite target-to-frame alignment (clergy~$\to$~abuse; institutions~$\to$~political/economic; identity labels~$\to$~rhetorical/moral) has the character of a genre convention --- a set of socially shared expectations about what kind of attack is appropriate for what kind of target.

\subsection{The Deep Social Logic Behind Language Community Divergence}

While the target-frame alignment holds cross-lingually, the moral architecture underlying it differs fundamentally between English and East Asian communities --- a divergence that becomes visible only through the RQ3 analysis and that requires a sociological rather than a purely linguistic explanation.

English-language anti-Christian discourse is negatively aligned across all five classical Moral Foundation axes simultaneously: Harm ($-0.074$), Fairness ($-0.063$), Loyalty ($-0.050$), Authority ($-0.081$), and Sanctity ($-0.065$). This broad negative sweep constitutes what may be described as \emph{full-spectrum moral condemnation}: the English-language speaker invokes Christianity's failure not on one ethical dimension but on all of them at once, arguing that the religion causes harm, enables injustice, demands blind loyalty, corrupts authority, and degrades human dignity. One possible explanation for this pattern lies in the cultural position of Christianity in Anglophone contexts: as the historical majority religion of Western societies, it is subject to an elevated standard of moral accountability. When institutions that have claimed moral authority are perceived to fail, the accusatory discourse that follows tends to indict them across the full range of moral standards they professed to uphold.

Chinese and Japanese communities present a structurally different profile. On the five classical Moral Foundation axes, both languages score positively --- Chinese: Harm ($+0.177$), Fairness ($+0.182$); Japanese: Harm ($+0.069$), Fairness ($+0.085$) --- but flip sharply negative on Symbolic Threat (ZH: $-0.047$; JA: $-0.045$) and Realistic Threat (ZH: $-0.012$; JA: $-0.018$). This pattern is consistent with the logic of \emph{identity-threat framing} rather than moral condemnation: the complaint is not that religion has violated shared ethical standards but that it poses a danger to cultural belonging, social cohesion, or group integrity. In this framing, Christianity figures not as a moral agent that has done something wrong, but as a foreign presence whose intrusion into the social fabric is itself the problem.

The divergence has a plausible socio-historical grounding. Chinese internet discourse about Christianity operates in a context of strong national secularism and a state-endorsed narrative that frames Western religion as incompatible with Chinese cultural tradition --- what 神棍 (\textit{religious fraudster}) encoding makes concrete: the clergy is attacked not only as morally deficient but as deliberately foreign and deceptive. Japanese anti-Christian discourse is further shaped by the specific political crisis of the Abe--Unification Church scandal, in which a South Korean-origin organisation's infiltration of domestic politics became the dominant frame; 韓国~(\textit{South Korea}) itself appears as a target in the Japanese top-20, encoding the Church's foreign origins as part of the critique. In this context, threat to national identity and political sovereignty --- Symbolic and Realistic Threat --- is the morally salient charge, while Western-style ethical condemnation of sexual misconduct or institutional corruption remains secondary.

\subsection{Integrating RQ1, RQ2, and RQ3: A Coherent Account}

When the three research questions are read together, a single coherent account of cross-lingual anti-Christian hate speech emerges. All three communities share a transnational core: Jesus/耶稣/イエス・キリスト and Catholic identity labels appear across languages, and the cross-linguistically balanced topics (the Conservative--Liberal Church rift, LGBTQ+ and Catholic teaching, the Science--Faith debate) confirm that some discourse themes transcend community-specific contexts. However, this shared surface conceals three structurally distinct modes of hostility.

The English-speaking community constructs its attacks through broad moral condemnation, targeting theological figures, clergy, and political actors alike with a rhetoric of ethical failure that mobilises all five moral foundations. The Chinese-speaking community concentrates its hostility on the clerical hierarchy and institutional Church, deploying the Sexual Abuse Frame and Institutional Rot as primary attack vectors and encoding the target in a register of deception and exploitation --- with 神棍 as the sharpest lexical crystallisation of this logic. The Japanese community's anti-Christian discourse is, in its dominant form, a discourse about domestic political corruption rather than theological controversy: the Church appears primarily as a vehicle for political capture and financial predation, and the moral charge is threat to national order and social safety rather than universal ethical failure.

This three-way divergence is not an artefact of topic composition or corpus sampling differences. The cross-language variation in moral framing persists even when the topic is held constant (Section~\ref{sec:rq3_topic_lang}): within the same topic, English texts consistently score negative on the classical Moral Foundation axes while Chinese and Japanese texts in the same topic score negative primarily on the Threat axes. The moral vocabulary difference is therefore a structural feature of each community's approach to anti-religious discourse, rooted in the social and political contexts that make different types of indictment locally resonant.

% ------------------------------------------------------------------
\section{Validity of the Who--How--Why Computational Framework}
\label{sec:discussion-pipeline}
% ------------------------------------------------------------------

Beyond the substantive findings, a core claim of this thesis is methodological: that the who--how--why pipeline offers a principled way to move computational hate speech research beyond binary detection toward structured, cross-linguistically comparable sociological interpretation. This section evaluates the scientific rationale for the two most distinctive architectural choices.

\subsection{Unsupervised Classification as the Architectural Core}

The decision to use BERTopic as the thematic scaffold for the entire pipeline --- rather than a predefined topic taxonomy --- was a deliberate epistemological choice, and one that requires explicit justification. Any predefined classification scheme applied prior to analysis would embed assumptions about which themes are salient in anti-Christian discourse. These assumptions would almost inevitably derive from English-language research traditions, where clerical sexual abuse and culture-war politics are the dominant documented phenomena. Imposing such a scheme on Chinese and Japanese corpora would risk \emph{failing to find} phenomena that are locally salient but absent from the prior literature --- precisely the Unification Church political crisis that constitutes the overwhelming majority of Japanese discourse, and the civilisational incompatibility framing characteristic of Chinese online nationalism.

By allowing BERTopic to induce the thematic structure directly from the corpus, the pipeline avoids this English-centric imposition. The empirical results validate the approach. Among the 47 extracted topics, the centroid alignment procedure successfully identified both ends of the cross-linguality spectrum: Topic~0 (Unification Church / Abe scandal) is 98.8\% Japanese, reflecting a discourse that has no structural equivalent in English or Chinese; Topic~43 (missionary and anti-Western sentiment) is 95.5\% Chinese; and Topic~16 is 91.8\% English. These monolingual topics are monolingual not because the alignment failed, but because the phenomena they describe are genuinely culture-specific. At the same time, the Science--Faith debate (Topic~7), the LGBTQ+ and Catholic Teaching topic, and the Conservative--Liberal Church rift each show near-perfectly balanced language distributions (entropy approaching 1.10), providing direct validation that the centroid alignment successfully brought Chinese, Japanese, and English documents on shared themes into the same topical cluster. No predefined taxonomy could have simultaneously captured this range of cultural specificity and cross-lingual consensus --- it emerges from the data itself.

\subsection{The Three-Layer Target Extraction for Zero-Subject Languages}

A structural challenge that the existing multilingual hate speech literature tends to understate is the pervasiveness of \emph{zero-subject constructions} in Chinese and Japanese. In contexts where the referent is recoverable from prior discourse or common ground, both languages routinely omit the grammatical subject and, frequently, the object as well. A sentence such as 性侵了（被害者）and covered it up has no parseable subject--verb--object triple in the surface form; standard dependency-based Named Entity Recognition, trained on structurally complete English sentences, produces a null result.

The three-layer extraction architecture --- spaCy NER with domain gazetteer enrichment followed by Gemini API fallback for structurally incomplete sentences --- was designed specifically to address this limitation. The NER and gazetteer layers handle the structurally complete cases efficiently; the LLM fallback recovers targets from the sentences that Stage~1 and Stage~2 cannot resolve, where context-dependent reference makes the attack object recoverable to a language model but not to a dependency parser. The result is a target extraction strategy that is adapted to the specific grammatical properties of the languages under analysis, rather than assuming that English-trained syntactic tools will generalise. The high entity counts recovered for Chinese (神父: 164 mentions; 教会: 306) and the precision of the Japanese target list (統一教会: 473; 安倍晋三: 92) suggest that the architecture is functioning as intended. This engineering decision is not merely a technical workaround; it reflects a principled commitment to building infrastructure that treats the linguistic properties of East Asian languages as first-class design constraints rather than as inconveniences to be patched around.

% ------------------------------------------------------------------
\section{Error Analysis and Technical Constraints}
\label{sec:discussion-errors}
% ------------------------------------------------------------------

Scientific credibility requires honest accounting of where the pipeline underperforms, introduces noise, or produces conclusions that are weaker than their surface form suggests. Three categories of limitation deserve explicit analysis: problems introduced by large language model usage at different pipeline stages; constraints inherent to the dictionary-based FrameAxis moral projection; and distortions in the training data introduced by back-translation augmentation.

\subsection{Three Modes of LLM Failure Across the Pipeline}

Large language models appear at three distinct points in the pipeline --- BERTopic topic labelling, RQ2 frame classification, and corpus augmentation --- and they introduce qualitatively different failure modes at each point.

\textbf{Under-determinism in BERTopic labelling.} The OpenAI API calls used to assign human-readable labels to raw BERTopic clusters are inherently non-reproducible: temperature-sampled generation means that rerunning the same prompt on the same cluster summary will produce a different label. This creates a tension between the deterministic character of scientific reporting and the stochastic character of the output. The topic labels used throughout this thesis (``Unification Church / Abe scandal'', ``LGBTQ+ and Catholic Teaching'', etc.) reflect one specific sampling run; a different run might produce labels of different granularity or emphasis, affecting the qualitative interpretation of the thematic landscape. The downstream quantitative analyses --- entity counts, frame distributions, moral bias scores --- are unaffected by this labelling variance, since they are computed on documents rather than topic labels; but the interpretive narrative built around those labels carries an uncertainty that standard reproducibility procedures cannot eliminate.

\textbf{Over-determinism in RQ2 frame classification.} The frame taxonomy applied in Chapter~\ref{chap:how} was defined prior to LLM classification: ten fixed categories were provided in the prompt, and the model was instructed to assign each predicate to the single most fitting category. While this design ensures cross-document comparability, it forecloses the possibility that the LLM might identify attack strategies that the predefined taxonomy does not anticipate. In pilot testing, the model rarely used its allowed freedom to flag ambiguous cases --- it consistently converged on one of the ten categories, even for predicates that did not fit cleanly. The result is that the frame distribution in Chapter~\ref{chap:how} reflects the coverage of the taxonomy rather than the full diversity of attack strategies in the corpus. If there exist rhetorical strategies in Chinese or Japanese anti-Christian discourse that have no analogue in the English-derived taxonomy from which the categories were partly constructed, those strategies are systematically suppressed rather than surfaced. A future iteration of the pipeline could address this by using an open-ended classification step first, followed by human-reviewed category consolidation, to let novel frame types emerge inductively.

\textbf{Semantic collapse in LLM-based corpus augmentation.} The back-translation augmentation chain used to expand the Chinese and Japanese hate speech training sets operates through a sequence of translation steps (English~$\to$~pivot languages~$\to$~target language). At each translation step, the model tends to normalise toward the most statistically frequent surface realisation of the underlying semantic content --- what has been described as semantic collapse in the multilingual generation literature. The practical consequence is that the augmented training data is likely to underrepresent the most culturally specific and rhetorically coded forms of anti-Christian hate speech: the abbreviations, wordplay, and community-specific coded terms that characterise 5ch (Japanese) and Tieba (Chinese) discourse are unlikely to survive a translation chain that passes through English, Korean, French, and German. The synthetic portion of the corpus (approximately 46\% of Chinese documents, 39\% of Japanese) may therefore produce a training signal that underweights the most challenging implicit hate speech cases precisely because the augmentation process strips them of their local flavour.

\subsection{FrameAxis Dictionary Constraints}

The FrameAxis moral projection framework relies on three lexical dictionaries: the Moral Foundations Dictionary for English, the Chinese Moral Foundations Dictionary \citep{cheng_c-mfd_2023}, and the Japanese Moral Foundations Dictionary \citep{matsuo2019development}. Two classes of limitation apply to the Chinese and Japanese dictionaries in particular.

First, \textbf{coverage gaps and polarity inference errors}. The C-MFD does not provide explicit polarity annotations; polarity was inferred via negation-prefix rules, which may misclassify compound terms where negation is semantically integrated rather than morphologically transparent. The J-MFD was developed for general Japanese text and may underrepresent the vocabulary of political and online discourse, where neologisms, abbreviations, and coded terminology proliferate. Both dictionaries are likely to assign incorrect or no scores to the most community-specific vocabulary --- exactly the portion of the corpus that carries the highest information value for understanding culturally specific hate speech. The specific axis scores reported in Chapter~\ref{chap:why} should therefore be interpreted as approximate indicators of moral orientation rather than precise measurements of moral activation.

Second, \textbf{interpretive ambiguity from high inter-axis correlation}. The five classical Moral Foundation axes in this corpus correlate at $r \approx 0.8$ --- substantially higher than the differentiated foundation profiles reported in survey-based MFT research \citep{graham2009liberals}, where liberals and conservatives show reliably distinct foundation profiles. This near-collapse of the five-dimensional foundation space into a single condemnation dimension creates genuine interpretive ambiguity: it is unclear whether hate speech as a communicative genre genuinely favours undifferentiated moral delegitimation (a substantive finding about hate speech's communicative logic), or whether the FrameAxis projection lacks the lexical resolution to distinguish foundations in online registers (a methodological artefact). The PCA analysis in Chapter~\ref{chap:why} --- which shows that PC1 alone explains 63.7\% of total variance --- provides a useful operational response to this ambiguity by focusing interpretive weight on the aggregate attack intensity and threat-versus-condemnation distinction captured by PC2, rather than claiming fine-grained precision at the individual-foundation level.

\subsection{Detection Model Errors and the Implicit Hate Speech Gap}

The fine-tuned hate speech classifiers used in corpus construction achieved Hate F1 scores of 0.635 (Chinese) and 0.579 (Japanese) at their best-performing epochs, representing substantial improvements over their untuned baselines (0.172 and 0.294 respectively) but still leaving substantial error rates. Error analysis of the misclassified cases reveals a systematic pattern: the classifier tends to misclassify \emph{implicit} hate speech --- posts that convey hostility through irony, cultural allusion, and coded vocabulary without using explicit slurs or overt threat language. For Chinese, this includes posts that mock Christianity through 神棍 framing or 毒瘤 (\textit{malignant tumour}) metaphor without explicit hate markers; for Japanese, it includes posts in the Unification Church discourse that treat political critique as a vehicle for religious hostility in ways that require cultural context to decode.

This failure mode has a positive methodological implication: it provides direct empirical support for the \emph{necessity} of the unsupervised BERTopic approach as the pipeline's architectural core. A supervised classifier trained to detect explicit hate speech will systematically exclude the implicit cases; an unsupervised topic model that clusters documents by semantic content rather than by hate-speech label will include implicit and explicit instances together, provided they cluster on shared topical content. The 2,530 documents assigned to Topic~$-1$ (noise) by BERTopic represent the cases that do not cluster cleanly --- but the 47 identified topics contain a mix of explicit and implicit hostility, with the implicit cases recoverable through the subsequent target extraction and frame analysis stages. The limitations of the supervised detection layer thus reinforce rather than undermine the case for the unsupervised architecture that follows it.

% ------------------------------------------------------------------
\section*{Chapter Summary}
% ------------------------------------------------------------------

The three sections of this chapter converge on a unified assessment of what the who--how--why pipeline has achieved and where it remains limited. The empirical findings reveal a structured, cross-linguistically comparable account of anti-Christian hate speech: target social category determines rhetorical frame selection in a genre-like fashion; English discourse operates through full-spectrum moral condemnation while East Asian discourse mobilises threat-based framing; and Japanese anti-religious discourse is, in its dominant form, a domestically-specific political discourse rather than theological hostility. The computational framework that produced these findings demonstrates the value of unsupervised topic modelling as an inductive alternative to English-centric predefined taxonomies, and of the three-layer extraction architecture as a principled response to the zero-subject structural challenge of CJK text. At the same time, the pipeline faces real constraints: LLM non-reproducibility in topic labelling, frame taxonomy rigidity that may suppress novel attack strategies, semantic collapse in augmented training data, dictionary gaps that limit the precision of moral axis scores, and systematic under-detection of implicit hate speech. These limitations affect the precision of specific quantitative estimates more than they affect the structural patterns themselves --- the cross-lingual divergences documented across all three research questions are robust to individual measurement errors because they emerge from large aggregate distributions rather than individual classifications. The following chapter situates these findings within the broader scholarly literature and articulates the theoretical and methodological contributions this thesis makes to computational hate speech research.

\subfilebibliography
\end{document}
